% ======================================================================
%  APPENDIX D: Engineering Principles Poster — tek sayfa
% ======================================================================
\clearpage
\thispagestyle{empty}
\markboth{Engineering Principles Poster}{Engineering Principles Poster}
\addcontentsline{toc}{chapter}{Appendix D · Principles Poster}

\begin{tikzpicture}[remember picture, overlay]
  \fill[accentdk] (current page.north west) rectangle
    ([yshift=-38mm]current page.north east);
  \node[anchor=north west, text=white, font=\sffamily\bfseries, inner sep=0pt]
    at ([shift={(19mm,-15mm)}]current page.north west)
    {\fontsize{12}{14}\selectfont 25 ENGINEERING PRINCIPLES};
  \node[anchor=north west, text=white, font=\sffamily\bfseries, inner sep=0pt]
    at ([shift={(19mm,-22mm)}]current page.north west)
    {\fontsize{22}{26}\selectfont Concepts That Matter};
  \node[anchor=south east, text=softgray, font=\sffamily\scriptsize, inner sep=0pt]
    at ([shift={(-17mm,16mm)}]current page.south east)
    {Araçlar eskir · ilkeler kalır};
\end{tikzpicture}

\vspace*{32mm}

% poster ilke öğesi: numara dairesi + kısa ifade
\newcommand{\pp}[2]{%
  \par\vspace{4.5pt}\noindent\hangindent=8mm\hangafter=0
  \makebox[7mm][l]{\tikz[baseline=-0.6ex]\node[fill=accent, text=white, circle,
    inner sep=0pt, minimum size=5mm, font=\sffamily\bfseries\scriptsize]{#1};}%
  \small #2\par}

\noindent
\begin{minipage}[t]{0.47\linewidth}
\pp{1}{AI'ı programlamazsın; brief verir, kısıtlar, ölçersin.}
\pp{2}{Prompt koddur: versiyonla ve test et.}
\pp{3}{Context bir bütçedir; her token bir tercih.}
\pp{4}{Anlamı, anlamı temsil eden uzayda ara.}
\pp{5}{Fine-tuning davranış içindir, bilgi için değil.}
\pp{6}{Modelin işi karar; eylemi koda yaptır.}
\pp{7}{Retrieval bir bilgi mimarisi sorunudur.}
\pp{8}{Bir kez yaz, her yerde kullan; sınırı sunucuda.}
\pp{9}{Hatırlama, doğru anda geri getirmektir.}
\pp{10}{Düşünmenin bedeli var; zorlukla orantılı öde.}
\pp{11}{Kontrol edebildiğin her şeyi workflow yap.}
\pp{12}{Adımları biliyorsan onları sen yaz.}
\pp{13}{Ölçemediğin kaliteyi iyileştiremezsin.}
\end{minipage}\hfill
\begin{minipage}[t]{0.47\linewidth}
\pp{14}{Göremediğin sistemi yönetemezsin.}
\pp{15}{Güvenliği prompt'a yazma; korkuluk kur.}
\pp{16}{En basitle başla; karmaşıklığı kanıt eklesin.}
\pp{17}{Daha büyük model çoğu sorunun cevabı değil.}
\pp{18}{Her başarısızlık yarının eval örneğidir.}
\pp{19}{``Bir kez çalıştı'' kanıt değildir.}
\pp{20}{Maliyet ve gecikme birer özelliktir.}
\pp{21}{Çıktıya körlemesine güvenme; doğrula.}
\pp{22}{Çoğu kalite sorunu bir girdi sorunudur.}
\pp{23}{Yüksek riskli kararda insanı döngüde tut.}
\pp{24}{Yapı sürprizi azaltır; yapılandır.}
\pp{25}{Bugünün modeline değil, değişime tasarla.}
\end{minipage}

\vfill
\begin{center}
\textcolor{accent}{\rule{30mm}{2pt}}\\[4pt]
{\sffamily\bfseries\small\color{accentdk}AI Engineering: Concepts That Matter}\\[1pt]
{\color{softgray}\itshape\footnotesize Kavramları değil, arkalarındaki düşünme biçimini öğren.}
\end{center}
\vspace{6mm}
